\documentstyle[% 


righttag% 


,11pt% 


,amssymb% 


,russian%               remove in Eng. 


,izvru%                 Set izven% in Eng. 


]{amsart} 


 


%  


\newcommand{\down}[2]{\underset{#1}{#2}{}} 


\newcommand{\sign}{\operatorname{sign}} 


\newcommand{\supp}{\operatorname{supp}} 


\newcommand{\tts}[1]{} 


 


%\hoffset=-15mm%                  For Eng. 


 


\setcounter{page}{3} 


\god{2000} 


\nomer{9~(460)} %                        Remove in Eng. 


%\nomer{Vol.\,44, No.\,9}%               Set in Eng. 


%\str{\pageref{begin}--\pageref{end}}%  Set in Eng. 


\udk{517.12.252} 


 


\author{\it Н.А.\,Ерзакова} 


% NB:   ^^ remove in Eng. 


\title{Меры некомпактности в исследовании неравенств} 


 


\date{} 


% \pagestyle{myheadings}%         Set in Eng. 


\pagestyle{plain} %              Remove in Eng. 


\begin{document} 


% \thispagestyle{plain}%          Set in Eng. 


\maketitle 


\label{begin} 


 


 


\section*{Введение} 


 


В ([1], с.\,74) рассматриваются банаховы пространства $E_0$, $E$, $E_1$, 


причем  


$E_0\subset E\subset E_1$, и вложение $E_0$ в $E$ компактно. 


Тогда для каждого $\varepsilon>0$ найдется такая постоянная 


$c_\varepsilon$, что для всех $v\in E_0$ 


\begin{equation} 


\|v\|_E\le\varepsilon\|v\|_{E_0}+ 


c_\varepsilon\|v\|_{E_1}. 


\end{equation} 


В ([2], с.\,154) изложен результат Ю.А.\,Дубинского, в котором 


$E_0$ заменяется множеством $\Im$, снабженным функцией $M:\Im\to R^+$, 


причем $\Im\subset E\subset E_1$, $M(v)\ge0$, 


\begin{equation} 


M(\lambda v)=|\lambda|M(v) 


\end{equation} 


для всех $v\in\Im$, множество $\Im_1$ относительно компактно в $E$, где 


\begin{equation} 


\Im_k=\{v\in\Im:M(v)\le k\} 


\end{equation} 


для всех $k\ge0$. 


При этих условиях для каждого $\varepsilon>0$ найдется такая постоянная 


$c_\varepsilon$, что 


\begin{equation} 


\|u-v\|_E\le\varepsilon(M(u)+M(v))+ 


c_\varepsilon\|u-v\|_{E_1} 


\end{equation} 


для всех $u$, $v$ из $\Im$. 


 


\begin{note}% 


[{\series{m}\shape{n}\selectfont [2], с.\,154}] Пусть $\Im=E_0$, $M$ --- 


норма в $E_0$. 


Тогда из неравенства (4) 


получается неравенство (1).  


\end{note} 


 


Введем величины 


$$ 


\tau(U)=\sup_{v\in U}M(v),\quad 


\wt\tau(U)=\inf_{v\in U}M(v). 


$$ 


 


\begin{defn} 


Множество $U\subset\Im$ назовем $M$-ограниченным, если $\tau(U)<\infty$. 


\end{defn} 


 


\begin{defn} 


Будем говорить, что множество $U\subset\Im$ обладает $M$-свойством, если  


оно $M$-ограничено и $\wt\tau(U)>0$. 


\end{defn} 


 


\begin{defn} 


Оператор (вообще говоря нелинейный) $A:\Im\to E$ назовем 


$M$-огра\-ни\-ченн\-ым, если любое $M$-ограниченное множество $U\subset\Im$ 


переводится оператором  


$A$ в ограниченное подмножество пространства $E$. 


\end{defn} 


 


Не предполагая (2), а также вложение $\Im_k$ компактным в $E$, 


будем исследовать условия  


справедливости утверждения: для каждого $\varepsilon>0$ 


найдется такая постоянная $c_\varepsilon$, что 


\begin{equation} 


\|A u-A v\|_E\le\varepsilon(M(u)-M(v))+c_\varepsilon 


||u-v\|_{E_1} 


\end{equation} 


для всех $u$, $v$ из произвольного множества $U\subset\Im$, обладающего 


$M$-свойством. 


 


\begin{note} 


В неравенстве (5) для каждого $\varepsilon>0$ постоянная $c_\varepsilon$, 


вообще говоря, 


зависит еще и от $U$. 


\end{note} 


 


Если $A$ --- линейный оператор, выполнено (2), кроме того, 


\begin{equation} 


M(v)>0\quad (v\ne0), 


\end{equation} 


то справедливость неравенства (5) для любого множества $U\subset\Im$, 


обладающего 


$M$-свойством, влечет его справедливость для всех $u$, $v$ 


из $\Im$ с одной и той же постоянной  


$c_\varepsilon$. Поскольку в ([2], с.\,154) в силу предположения (2)  


и относительной компактности $\Im_1$ 


имеет место (6), неравенства (1) и (4)  


являются частными случаями неравенства (5).  


 


Отличительной чертой исследований, проводимых в данной работе, является 


применение мер некомпактности. 


 


\begin{defn}% 


[{\series{m}\shape{n}\selectfont [3], с.\,20}] Пусть $E$ --- банахово 


пространство, а $U$ --- ограниченное  


подмножество $E$. Мерой некомпактности $\beta_E(U)=\beta(U)$ множества 


$U$ называется  


супремум тех $r>0$, для которых в $U$ 


существует такая бесконечная последовательность  


$\{u_n\}$, что $\|u_n-u_m\|\ge r$ для всех $n\ne m$. 


\end{defn} 


 


Мерой некомпактности Хаусдорфа $\chi_E(U)=\chi(U)$ множества $U$ 


называется инфимум  


всех $\varepsilon>0$, при которых $U$ имеет в $E$ конечную 


$\varepsilon$-сеть.  


 


Итак, в ([1], с.\,71; [2], с.\,154) 


приведены достаточные условия справедливости 


неравенств (1) и (4), цель же данной работы --- 


получить необходимые и достаточные  


условия справедливости обобщающего неравенства (5). 


 


\section{Основные результаты} 


 


\begin{thm} 


Пусть $A$ --- $M$-ограниченный оператор из $\Im$ в $E$. 


Неравенство $(5)$ справедливо, если и только если выполнено 


следующее условие$:$ 


 


если последовательности $\{u_n\}$, $\{v_n\}$ обладают $M$-свойством и, 


кроме того, 


\begin{equation} 


\lim_{n\to\infty}\|u_n-v_n\|_{E_1}=0, 


\end{equation} 


то 


\begin{equation} 


\down{n\to\infty}{\varliminf} 


\|A u_n-A v_n\|_E=0.


\end{equation} 


\end{thm} 


 


\begin{pf} 


Предположим, что условие теоремы выполнено. Для доказательства  


справедливости неравенства (5) предположим противное: 


неравенство (5) не выполнено. Последнее означает, что найдутся  


множество $U$, обладающее $M$-свойством, $\varepsilon>0$, 


последовательности  


элементов $\{u_n\}$, $\{v_n\}$ ($u_n\in U$, $v_n\in U$ для каждого $n$) 


и чисел $\{c_n\}$ ($c_n\to\infty$),  


для которых справедливо неравенство 


\begin{equation} 


\|A u_n-A v_n\|_E>\varepsilon(M(u_n)-M(v_n))+ 


c_n\|u_n-v_n\|_{E_1} 


\end{equation} 


при всех $n$. 


 


По предположению оператор $A$ является $M$-ограниченным. 


Так как $\{u_n\}$, $\{v_n\}$ обладают $M$-свойством и, следовательно, 


$M$-ограничены, то последовательности $\{A{}u_n\}$, $\{A{}v_n\}$ 


равномерно ограничены в $E$. Отсюда 


последовательность $\{A{}u_n-A{}v_n\}$ также равномерно  


ограничена в $E$. В силу неравенства (9) и предположения о 


том, что $c_n\to\infty$, последнее  


возможно, если $\|u_n-v_n\|_{E_1}\to0$. Поэтому имеет место (8). 


Полученное противоречие с (9) завершает доказательство. 


 


Предполагая теперь неравенство (5) выполненным,  


докажем справедливость условия теоремы. 


Пусть последовательности $\{u_n\}$, $\{v_n\}$ обладают $M$-свойством. 


Это означает, что 


найдутся числа $0<r\le R$, для которых справедливо неравенство $r\le 


M(u_n)+M(v_n)\le R$ 


при всех $n$. И пусть, кроме того, имеет место (7). 


Покажем, что справедливо (8). 


 


С целью построения нужной подпоследовательности  


для произвольного $\varepsilon>0$ выберем  


из $\{A{}u_n-A{}v_n\}$ элемент 


$A{}u_{n_\varepsilon}-A{}v_{n_\varepsilon}$ с 


номером $n_\varepsilon$, удовлетворяющим неравенству 


$$ 


\|u_{n_\varepsilon}-v_{n_\varepsilon}\|_{E_1}\le 


\varepsilon/c_\varepsilon, 


$$ 


где $c_\varepsilon$ --- постоянная из (5). Из неравенства (5) получим 


$$ 


\|A u_{n_\varepsilon}-A v_{n_\varepsilon}\|_E\le 


\varepsilon R+\varepsilon, 


$$ 


откуда следует справедливость условия теоремы. 


\end{pf} 


 


В последующих результатах этого параграфа не предполагается, что  


множество $U\subset\Im$ обладает $M$-свойством. 


 


\begin{thm} 


Пусть задан $M$-ограниченный оператор $A:\Im\to E$. Пусть для 


некоторого $\varepsilon>0$ 


найдется такая постоянная $c_\varepsilon$, что неравенство $(5)$ 


справедливо для всех $u$, $v$ 


из $\Im$. Тогда для этих же постоянных $\varepsilon$ и $c_\varepsilon$ 


имеет место неравенство 


\begin{equation} 


\beta_E(A U)\le2\varepsilon\tau(U)+ 


c_\varepsilon\beta_{E_1}(U) 


\end{equation} 


для каждого множества $U\subseteq\Im$. 


\end{thm} 


 


\begin{pf} 


Итак, пусть неравенство (5) справедливо для всех $u$, $v$ из $\Im$. 


Пусть $U$ --- произвольное подмножество $\Im$. Если $\beta_E(A{}U)=0$, 


то утверждение очевидно.  


Пусть $\beta_E(A{}U)>0$. По определению меры некомпактности $\beta_E$ 


для произвольного $0<\delta<\beta_E(A{}U)$ 


найдется такая последовательность $\{A{}u_n\}\subset A{}U$, что для 


всех $n\ne m$ 


$$ 


\beta_E(A U)-\delta\le\|A u_n-A v_m\|_E. 


$$ 


Из определения $\beta$ нетрудно вывести, что в $\{u_n\}$ можно выбрать 


элементы $\wt u_n$, $\wt u_m$, 


удовлетворяющие неравенству 


$$ 


\|\wt u_n-\wt u_m\|_{E_1}\le 


\beta_{E_1}(U)+\delta. 


$$ 


По предположению для некоторых $\varepsilon$ и $c_\varepsilon$ 


справедливо  


неравенство (5). Поэтому, применяя  


неравенство (5), получим 


$$ 


\beta_E(A U)-\delta\le\|A\wt u_n-A\wt u_m\|_E\le 


\varepsilon(M(\wt u_n)+M(\wt u_m))+c_\varepsilon 


\|\wt u_n-\wt u_m\|_{E_1}\le2\varepsilon\tau(U) 


+c_\varepsilon(\beta_{E_1}(U)+\delta), 


$$ 


откуда в силу произвольности выбора $\delta$ следует (10). 


\end{pf} 


 


\begin{prop} 


Пусть для всех $\varepsilon>0$ найдется такая постоянная $c_\varepsilon$, 


что для каждого  


$U\subseteq\Im$ имеет место неравенство $(10)$. Тогда если 


$U\subseteq\Im$ является $M$-ограниченным и  


относительно компактным в $E_1$, то $A{}U$ также относительно компактно 


в $E$.  


\end{prop} 


 


\begin{pf} 


Допустим, что $U\subseteq\Im$ относительно компактно в $E_1$. Тогда в  


силу правильности меры некомпактности $\beta_E$ ([2], с.\,7) 


получим $\beta_{E_1}(U)=0$. Отсюда и из  


неравенства (10) справедливо неравенство 


$\beta_E(A{}U)\le\varepsilon\tau(U)$. Так как $\varepsilon>0$ 


произвольно, то  


$\beta_E(A{}U)=0$. Это означает, что $A{}U$ относительно компактно в 


$E$. 


\end{pf} 


 


\section{Приложения для линейного оператора} 


 


\begin{prop} 


Пусть множество $\Im_k$, определенное равенством $(3)$, относительно  


компактно в $E$ для всех $k\ge0$, $E=E_0$, $A$ --- тождественный 


оператор, тогда справедливо  


неравенство $(5)$ для $U$, обладающих $M$-свойством.  


\end{prop} 


 


\begin{pf} 


В предположениях доказываемого утверждения выполнено условие  


теоремы 1. Действительно, если последовательности $\{u_n\}$, $\{v_n\}$ 


обладают 


$M$-свойством и, кроме того, имеет место (7), 


то в силу относительной компактности последовательностей 


$\{u_n\}$, $\{v_n\}$ в $E$ имеем 


$$ 


\down{n\to\infty}{\varliminf}\|u_n-v_n\|_E=0. 


$$ 


Отсюда, учитывая эквивалентность условия теоремы 1 


неравенству (5), получаем справедливость (5). 


\end{pf} 


 


\begin{note} 


Принимая во внимание замечания 1, 2, приходим к выводу, что предложение 


2 является обобщением теоремы 16.4 из ([1], с.\,126) 


и леммы 12.1 из ([2], с.\,154), в  


которой предполагается (2). 


\end{note} 


 


\begin{thm} 


Пусть $A$ --- оператор вложения из $E_0$ в $E$, $\Im=E_0$, $M$ --- 


норма в $E_0$, и любое  


множество $U$, ограниченное в $E_0$ 


и относительно компактное в $E_1$, также относительно  


компактно в $E$. Тогда выполнено неравенство $(5)$ для всех $u$ и $v$. 


\end{thm} 


 


\begin{pf} 


Предполагая все предположения теоремы 3 выполненными, покажем 


справедливость условия теоремы 1. 


Пусть последовательность $\{u_n-v_n\}$ такова, что $u_n,v_n\in E_0$ для 


всех $n$, $0<r\le\|u_n\|_{E_0}+\|v_n\|_{E_0}\le R$ 


для некоторых $r$ и $R$ и, кроме того, имеет место (7). 


Но тогда, поскольку $\{u_n-v_n\}$ относительна компактна в $E_1$, 


по предположению доказываемой теоремы $\{u_n-v_n\}$ 


также относительно компактна в $E$ и для некоторой подпоследовательности 


$\{\wt u_n-\wt v_n\}$ найдется такой элемент $u_0\in E$, что 


$\lim\limits_{n\to\infty}\|(\wt u_n-\wt v_n)-u_0\|_E=0$. 


В силу вложения $E\subset E_1$ имеем также 


$\lim\limits_{n\to\infty}\|(\wt u_n-\wt v_n)-u_0\|_{E_1}=0$. 


Отсюда, учитывая свойство единственности предела, $u_0=0$, т.\,е. 


выполнено условие теоремы 1, которое эквивалентно неравенству 


(5). Принимая во внимание замечание 2,  


получим утверждение теоремы 3. 


\end{pf} 


 


\begin{cor} 


Пусть $A$ --- оператор вложения из $E_0$ в $E$, $\Im=E_0$, $M$ --- 


норма в $E_0$.  


Тогда неравенство (1) эквивалентно каждому из трех условий: 


\begin{itemize} 


\item[$1'$.] для любой ограниченной последовательности $\{u_n\}$ из 


$E_0$, удовлетворяющей равенству 


$\lim\limits_{n\to\infty}\|u_n\|_{E_1}=0$,


имеем также $\down{n\to\infty}{\varliminf}\|u_n\|_E=0$; 


\item[$2'$.] для каждого $\varepsilon>0$ найдется такая постоянная 


$c_\varepsilon$, что для каждого $U\subset E$, 


равномерно ограниченного в $E_0$, справедливо неравенство 


$\beta_E(U)\le\varepsilon\beta_{E_0}(U)+ 


c_\varepsilon\beta_{E_1}(U)$; 


\item[$3'$.] любое множество $U$, ограниченное в $E_0$ и относительно 


компактное в $E_1$, также  


относительно компактно в $E$. 


\end{itemize} 


\end{cor} 


 


{\bf Пример.} Покажем, что в ([1], с.\,71; [2], с.\,154) 


предположения относительно компактности не являются необходимыми. 


 


Рассмотрим на отрезке $(0,1)$ пространства Лебега $L_\infty$, $L_2$, 


$L_1$. 


Известно, что имеет место вложение $L_\infty\subset L_2\subset L_1$. 


Вложение $L_\infty$ в $L_2$ некомпактно.  


Действительно, ограниченная по норме в $L_\infty$ 


последовательность функций Радемахера  


$\{u_n\}$, где $u_n(t)=\sign\sin 2^{n-1}\pi t$, 


имеет в пространстве $L_2$ меру некомпактности $\beta$,  


равную, как легко видеть, $2^{1/2}$.  


 


Однако непосредственно проверяется справедливость  


утверждения: для каждого $\varepsilon>0$ 


найдется такая постоянная $c_\varepsilon$, что для всех $v\in L_\infty$ 


$$ 


\|v\|_{L_2}\le\varepsilon\|v\|_{L_\infty}+ 


c_\varepsilon\|v\|_{L_1}. 


$$ 


 


В завершение работы приведем приложения полученных  


результатов к оператору вложения пространств Соболева. 


 


Пусть $\Omega$ --- область в $\cal R^n$ с конечной мерой Лебега, т.\,е. 


$m(\Omega)<\infty$, $L^{l,p}(\Omega)=L^{l,p}$ 


($1\le p<\infty$) --- множество таких обобщенных функций на $\Omega$, 


что $\|\nabla_l u\|_{L_p}<\infty$, где 


$$ 


\|\nabla_l u\|_{L_p}= 


\bigg(\int_\Omega\sum_{|\alpha|=l}|D^\alpha u(x)|^p 


dx \bigg)^{1/p}. 


$$ 


Пусть $W^{k,p}(\Omega)=W^{k,p}$ --- пространство Соболева с нормой 


$$ 


\|u\|_{W^{k,p}(\Omega)}= 


\bigg(\sum^k_{l=0}\|\nabla_l u\|^p_{L_p}\bigg)^{1/p}. 


$$ 


Пусть $0\le k\le l<s$. В ([4], с.\,33) для ограниченной области 


$\Omega$ с кусочно-гладкой  


границей приводится неравенство Эрлинга--Ниренберга 


\begin{equation} 


\|u\|_{W^{l,p}}\le\varepsilon\|u\|_{W^{s,p}}+ 


c_\varepsilon\|u\|_{W^{k,p}}, 


\end{equation} 


справедливое для всех $\varepsilon>0$. 


 


В силу определения нормы в пространстве Соболева всегда имеет место 


вложение  


$W^{s,p}\subset W^{l,p}\subset W^{k,p}$ при $0\le   k\le l<s$. Однако 


если не предполагать границу области  


кусочно-гладкой, то, как следует из ([5], с.\,212), 


вложение $W^{s,p}$ в $W^{l,p}$ при $l<s$ может  


оказаться некомпактным. 


 


\begin{cor} 


Неравенство (11) справедливо тогда и только тогда, когда относительная  


компактность каждого ограниченного подмножества $W^{s,p}$ в $W^{k,p}$ 


влечет его относительную компактность в $W^{l,p}$. 


\end{cor} 


 


Для полноты освещения затронутого вопроса использования мер некомпактности  


в  


неравенствах докажем аналог неравенства Эрлинга--Ниренберга, справедливый 


также для некомпактного оператора вложения $W^{s,2}$ в $W^{l,2}$. 


 


Пусть $C$ --- пространство всех постоянных функций на $\Omega$ 


и пусть оператор вложения 


\begin{equation} 


I:L^{1,2}/C\to L_2/C 


\end{equation} 


ограничен. 


Пусть $C^{0,1}(\Omega)=C^{0,1}$ --- пространство функций, 


удовлетворяющих условию Липшица  


на любом компактном подмножестве множества $\Omega$. 


Для произвольного $T>0$ и $u\in L_2$ пусть 


$D(u,T)=\{x\in\Omega:|u(x)|>T\}$, а для  


произвольного $\varepsilon>0$ обозначим 


\begin{equation} 


U(\varepsilon)=\sup_{m(D)<\varepsilon}\sup 


\{\|U\|_{L_2}:u\in C^{0,1},\ \supp 


u\subseteq D,\ \|\nabla u\|_{L_2}=1\}. 


\end{equation} 


 


Из теоремы 4 работы [6] следует, что 


$$ 


\inf\{U(\varepsilon):\varepsilon>0\}=\chi(IS), 


$$ 


где $\chi(IS)$ --- мера некомпактности Хаусдорфа (определение 4) 


множества $\{u\in L^{1,2}:\|\nabla u\|_{L_2}=1\}$, 


вычисленная в фактор-пространстве $L_2/C$. 


Величина $\chi(IS)$ является характеристикой степени некомпактности 


оператора $I$, 


поскольку обращается в нуль, если и только если оператор (12) компактен.  


 


\begin{prop} 


Для любых $\varepsilon>0$ и $u\in L_2\cap L^{l+1,2}$ справедливо 


неравенство 


$$ 


\|u\|_{W^{l,2}}\le C(U(\varepsilon),n) 


\|\nabla_{l+1}u\|_{L_2}+c_\varepsilon\|u\|_{L_2}, 


$$ 


где $C(U(\varepsilon),n)$ зависит только от $U(\varepsilon)$ и $n$ и, в 


частности, 


$$ 


C(U(\varepsilon),n)=a U(\varepsilon)/ 


(1-a U(\varepsilon)), 


$$ 


если $a{}U(\varepsilon)<1$ $(a=n^{1/2}+1)$. 


\end{prop} 


 


\begin{pf} 


Так же как в теореме 4.8.2 из ([5], с.\,203), для произвольно выбранных  


$\varepsilon>0$ и $u\in L_2\cap L^{l+1,2}$ положим 


$T=\inf\{t:m(D(u,t))\le\varepsilon\}$. В силу теоремы 1.1.1 из ([5], 


с.\,16) 


без ограничения общности можем предполагать, что $u\in 


C^\infty(\Omega)$. Из (13) имеем 


$$ 


\|u\|_{L_2}\le\|\,|u|-T\|_{L_2(D(u,T))}+ 


\|\,|u|-T\|_{L_2(\Omega\setminus D(u,T))}+ 


T[m(\Omega)]^{1/2}\le U(\varepsilon) 


\|\nabla u\|_{L_2}+2T[m(\Omega)]^{1/2}.


$$ 


Обозначим через $\Omega_\varepsilon$ такую ограниченную подобласть


$\Omega$


с границей класса $C^{0,1}$, что


$m(\Omega\setminus\Omega_\varepsilon)<\varepsilon/2$. Так как


$m[D(u,T)]\ge\varepsilon$, то


$m[D(u,T)\cap\Omega_\varepsilon]\ge\varepsilon/2$.


Следовательно,


$\|u\|_{L_r(\Omega_\varepsilon)}\ge 2^{-1/r}T\varepsilon^{1/r}$


для произвольного $r\ge1$. Поэтому


$$


\|u\|_{L_2(\Omega)}\le U(\varepsilon)


\|\nabla u\|_{L_2(\Omega)}+2^{1-1/r}


[m(\Omega)]^{1/2}\|u\|_{L_r(\Omega)},


$$


где оператор вложения из $L^{1,2}(\Omega_\varepsilon)$ в 


$L_r(\Omega_\varepsilon)$ компактен. Дальнейшие рассуждения  


аналогичны рассуждениям, проведенным при доказательстве леммы 


4.10.2 из ([5], с.\,211).


\end{pf} 


 


\begin{note} 


Предложение 3 обобщает лемму 4.10.2 из ([5], с.\,211) 


на случай некомпактного оператора вложения (12). 


\end{note} 


 


Как показано в [6], 


для доказательства разрешимости задачи Неймана в некоторых  


случаях не нужно требовать справедливость (11) 


для всех $\varepsilon>0$. Достаточно, чтобы $\chi(IS)$ 


была меньше заданной величины. 


 


\begin{thebibliography}{9} 


 


\bibitem{1} 


Лионс Ж.-Л., Мадженес Э. {\em Неоднородные граничные задачи и их 


приложения}. 


-- М.: Мир, 1971. -- 371 с. 


\bibitem{2} 


Лионс Ж.-Л. {\em Некоторые методы решения нелинейных краевых задач}. -- 


М.: Мир, 1972. -- 588~с. 


\bibitem{3} 


{\em Меры некомпактности и уплотняющие операторы} / Под ред. 


P.P.\,Ахмерова, М.И.\,Каменского, А.С.\,Потапова,


Б.Н.\,Садовского и др. -- Новосибирск: Наука, 1986. -- 265 с. 


\bibitem{4} 


Березанский Ю.М. {\em Разложение по собственным функциям самосопряженных  


операторов}. -- Киев: Наук. думка, 1965. -- 798 с. 


\bibitem{5} 


Мазья В.Г. {\em Пространства С.Л.\,Соболева}. -- 


Л.: Изд-во Ленингр. ун-та, 


1985. -- 415 с. 


\bibitem{6} 


Yerzakova N.A. {\em On measures of non-compactness in regular spaces} // 


Zeitshcr. Anal. Anw. -- 1996. -- V.\,15. -- \No\,2. -- P.\,299--307. 


 


\end{thebibliography} 


 


\vskip 2cm 


 


\post{\it Хабаровский государственный}{\it Поступила} 


\post{\it технический университет}{22.01.1999} 


 


\label{end} 


\end{document} 


 


